% Píše: backend. Krátký český abstrakt celé práce.
Práce popisuje návrh a~realizaci simulovaného chytrého domku, který se ovládá z~tabletu
s~Androidem přes Bluetooth Low Energy. Řídicí jednotkou domku je Raspberry Pi 3B+ s~programem
v~Pythonu. Ten čte skutečná čidla (teplotu a~tlak, osvětlení a~pohyb), simuluje vnitřní teplotu
domu a~termostat s~topením a~chlazením, automaticky řídí barevné světlo podle denní doby,
osvětlení a~pohybu a~stav ukazuje na displeji a~kontrolkách. Navenek vystupuje jako periferie
BLE se standardní službou Environmental Sensing a~vlastní službou, která přenáší stav, nastavení
a~příkazy ve formátu JSON. Mobilní aplikace v~Kotlinu stav zobrazuje a~domek ovládá. Součástí
práce je popis zapojení, komunikačního rozhraní, obou programů a~jejich ověření automatickými
testy i~na skutečném hardwaru.
